\subsection{Problem Formulation}

We model redistricting as constrained graph partitioning. Census tracts form graph vertices with 2D weights $w_v = (p_v, m_v)$ representing total population and minority VAP. Edges connect geographically adjacent tracts. Goal: partition into $k$ districts minimizing edge cuts (compactness proxy) while satisfying:

\begin{align}
\text{Population balance:} \quad & \left|\frac{\sum_{v \in D_i} p_v}{\sum_v p_v} - \frac{1}{k}\right| \leq \epsilon_p && \forall i \label{eq:pop}\\
\text{Contiguity:} \quad & D_i \text{ forms connected subgraph} && \forall i \label{eq:contig}\\
\text{MM target:} \quad & \frac{\sum_{v \in D_i} m_v}{\sum_{v \in D_i} p_v} \geq 0.50 && \text{for } t \text{ districts} \label{eq:vra}
\end{align}

where $D_i$ is district $i$, $\epsilon_p = 0.005$ is population tolerance, and $t$ is the target number of MM districts (typically 2-5 depending on state).

\subsection{States and Targets}

We test five VRA-covered states with varying demographics and district counts:

\begin{table}[h]
\centering
\begin{tabular}{lcccc}
\toprule
State & Districts & Target MM & Minority \% & Tracts \\
\midrule
Alabama (AL) & 7 & 2 & 36.9\% & 1,437 \\
Georgia (GA) & 14 & 5 & 42.4\% & 2,796 \\
Louisiana (LA) & 6 & 2 & 41.6\% & 1,388 \\
Mississippi (MS) & 4 & 2 & 46.1\% & 878 \\
South Carolina (SC) & 7 & 3 & 35.1\% & 1,323 \\
\bottomrule
\end{tabular}
\caption{VRA test states with 2020 Census demographics}
\label{tab:states}
\end{table}

Target MM district counts reflect statutory requirements and recent legal precedent. Minority VAP includes Black and Hispanic populations, following standard VRA analysis \citep{katz2020}.

\subsection{Data}

\textbf{Geographic}: Census 2020 TIGER/Line shapefiles for tract boundaries and adjacency.

\textbf{Demographic}: Census 2020 redistricting files (PL 94-171) for total population and voting-age population (VAP) by race/ethnicity. \emph{Note on denominator}: The MM district threshold (50\% minority) and the minority tract identification threshold $\tau$ are applied against \textbf{total population} (not VAP) throughout this paper. Total population is the appropriate denominator for redistricting population-balance calculations under the Equal Protection Clause; VAP is the denominator courts examine for VRA dilution analysis. Where the background section references ``minority VAP'' in the context of Gingles preconditions, this refers to the legal standard; our algorithmic threshold uses total population as a practical proxy because PL 94-171 total population data is available at the tract level for all 50 states.

\textbf{Graph Construction}: Queen contiguity (shared boundary or vertex) with water-based connections for disconnected components. Final graphs: 878-2,796 vertices, average degree 6.1.

\subsection{Optimization Approaches}

We test four approaches combining two partitioning methods with two constraint modes:

\subsubsection{Method 1: Recursive Bisection}

Hierarchically split into districts by repeated binary partitioning \citep{dellaLibera2026recursive}. At each split of $n$ districts into $(n_L, n_R)$:

\begin{enumerate}
\item Calculate target weights for left/right parts based on desired MM distribution
\item Run METIS 2-way partition with target weights
\item Recurse on each part until all districts formed
\end{enumerate}

Tree structure matters: different binary trees yield different intermediate constraints. We test predetermined balanced trees (e.g., $7 \to [3,4] \to \ldots$).

\subsubsection{Method 2: N-Way Partitioning}

Partition directly into all $k$ districts simultaneously. Specify target weights for each final district:
\begin{itemize}
\item MM districts ($i = 1, \ldots, t$): Equal population, high minority share
\item Non-MM districts ($i = t+1, \ldots, k$): Equal population, remaining minority
\end{itemize}

N-way considers all edge cuts globally, avoiding intermediate constraint limitations of recursive bisection.

\subsubsection{Method 3: Adaptive Recursive Bisection}

At each split, evaluate both $(n_L, n_R)$ and $(n_R, n_L)$ orderings. Choose whichever achieves better minority concentration. Makes locally optimal decisions rather than following predetermined tree.

\subsubsection{Constraint Modes}

\textbf{Strict (tpwgts only)}: Specify target weights $w_{ij}$ for partition $i$, constraint $j$. METIS balances both population and minority constraints simultaneously with tolerance $\epsilon = 0.005$.

\textbf{Relaxed (tpwgts + ubvec)}: Additionally specify per-constraint imbalance: $\text{ubvec} = [1.005, 1000]$. Allows 1000$\times$ more imbalance in minority distribution, enabling concentration while maintaining population balance.

\subsection{Target Weight Calculation}

For $k$ districts with target $t$ MM districts at 60\% minority:

\begin{align}
\text{MM districts:} \quad w_{i,\text{pop}} &= \frac{1}{k}, \quad w_{i,\text{min}} = \frac{0.60/k}{m_{\text{state}}} && i \leq t\\
\text{Non-MM:} \quad w_{i,\text{pop}} &= \frac{1}{k}, \quad w_{i,\text{min}} = \frac{(m_{\text{state}} - 0.60t/k)/k_{\text{non}}}{m_{\text{state}}} && i > t
\end{align}

where $m_{\text{state}}$ is state-wide minority fraction and $k_{\text{non}} = k - t$ is number of non-MM districts.

\subsection{Method 4: Edge-Weighted Single-Objective Optimization}

Multi-constraint optimization (Methods 1-3) treats VRA compliance as a dual-objective problem, simultaneously balancing population and minority VAP. However, this approach proved insufficient for states with dispersed minority populations. We develop an alternative: \textbf{single-objective optimization with weighted edges}.

\subsubsection{Edge Weighting Strategy}

Rather than multi-constraint balancing, we modify the graph structure itself. For each edge $(u, v) \in E$ connecting tracts $u$ and $v$:

\begin{equation}
w_{edge}(u, v) = \begin{cases}
\alpha & \text{if } m_u \geq \tau \text{ and } m_v \geq \tau \\
1 & \text{otherwise}
\end{cases}
\end{equation}

where $m_u$ is the minority percentage in tract $u$, $\tau$ is a minority threshold (e.g., 0.45 or 0.50), and $\alpha$ is a weight factor (e.g., 5, 10, 100).

\textbf{Key insight}: Edges between high-minority tracts receive higher weights, making them "expensive" to cut. METIS minimizes total weighted edge cut, so it naturally preserves minority communities.

\subsubsection{Optimization Objective}

Unlike Methods 1-3 which use 2D vertex weights $w_v = (p_v, m_v)$, Method 4 uses:
\begin{itemize}
\item \textbf{Vertex weights}: 1D (population only): $w_v = p_v$
\item \textbf{Edge weights}: As defined above
\item \textbf{No target weights} (no tpwgts file)
\item \textbf{No constraint relaxation} (no ubvec parameter)
\end{itemize}

This is pure single-objective edge-cut minimization with a modified graph. METIS minimizes:
\begin{equation}
\text{EdgeCut} = \sum_{(u,v) \in E : D_u \neq D_v} w_{edge}(u, v)
\end{equation}

subject only to population balance constraint (\ref{eq:pop}).

\subsubsection{Ablation Study Design}

We systematically test edge weighting across parameter space:
\begin{itemize}
\item \textbf{Weight factors} $\alpha$: 1, 5, 10, 50, 100, 500, 1000
\item \textbf{Minority thresholds} $\tau$: 0.40, 0.45, 0.50, 0.55
\item \textbf{States}: All 5 VRA-covered (AL, GA, LA, MS, SC)
\end{itemize}

Total configurations: $7 \times 4 \times 5 = 140$ tests.

For each configuration, we measure:
\begin{enumerate}
\item MM districts achieved (primary metric)
\item Maximum minority concentration
\item Edge cut (compactness proxy): unweighted and weighted
\item Average internal edges per district
\end{enumerate}

\subsubsection{Comparison to Multi-Constraint}

Edge weighting differs fundamentally from Methods 1-3:

\begin{table}[h]
\centering
\small
\begin{tabular}{lcc}
\toprule
& Multi-Constraint & Edge-Weighted \\
\midrule
Vertex weights & 2D (pop + minority) & 1D (pop only) \\
Edge weights & Uniform (1.0) & Variable ($1$ or $\alpha$) \\
Optimization & Dual-objective & Single-objective \\
Target weights & Yes (tpwgts) & No \\
Constraint relaxation & Optional (ubvec) & Not applicable \\
Complexity & High & Low \\
\bottomrule
\end{tabular}
\caption{Multi-constraint vs edge-weighted optimization}
\label{tab:method_comparison}
\end{table}

\textbf{Hypothesis}: Edge weighting directly encodes the VRA goal ("don't split minority communities") into graph structure, potentially achieving better results than indirect multi-constraint balancing.

\subsection{Production Algorithm: Adaptive Weight Formula}
\label{sec:adaptive_formula}

The ablation study above tests fixed weight factors to characterize the parameter space. The deployed V4 pipeline uses an \textbf{adaptive formula} that sets $\alpha$ automatically based on each state's minority tract density:

\begin{equation}
\alpha = \max\!\left(3.0,\; 10.0 \times \bigl(1.0 - 0.7 \times f_{\text{minority}}\bigr)\right)
\label{eq:adaptive_boost}
\end{equation}

where $f_{\text{minority}}$ is the fraction of tracts in the state with minority percentage $\geq \tau$ (the minority threshold, default $\tau = 0.40$).

\textbf{Rationale}: States with densely clustered minority tracts (high $f_{\text{minority}}$) need less boosting---the geography already concentrates minority tracts together. States with sparse minority tracts (low $f_{\text{minority}}$) need stronger signals to overcome dispersion. The formula scales from 10$\times$ (no minority tracts, $f = 0$) down to a floor of 3$\times$ (saturated states), with the $0.7$ coefficient tuned to keep $\alpha$ in the empirically effective 5$\times$--10$\times$ range for the covered Southern states.

\textbf{Effective weights for key states} (2020 Census, $\tau = 0.40$):
\begin{center}
\begin{tabular}{lccc}
\toprule
State & $f_{\text{minority}}$ & $\alpha$ (adaptive) & Closest fixed config \\
\midrule
Alabama      & $\sim$0.22 & $\approx 8.5\times$ & 10$\times$ \\
Georgia      & $\sim$0.34 & $\approx 7.6\times$ & 10$\times$ \\
Louisiana    & $\sim$0.30 & $\approx 7.9\times$ & 10$\times$ \\
Mississippi  & $\sim$0.47 & $\approx 6.7\times$ & 5$\times$ \\
South Carolina & $\sim$0.16 & $\approx 8.9\times$ & 10$\times$ \\
California   & $\sim$0.63 & $\approx 5.6\times$ & 5$\times$ \\
\bottomrule
\end{tabular}
\end{center}

\textbf{Relationship to ablation}: The adaptive formula maps each state to an effective $\alpha$ within the 5$\times$--10$\times$ range that the ablation identifies as optimal (Section~\ref{sec:ablation_results}). Results for the ablation's fixed-weight configurations are therefore directly comparable to V4 production outputs for all covered states except those near the 3$\times$ floor (high-diversity states such as California, Hawaii).

\textbf{V4 production outcomes} (verified 2026-04-24, current pipeline):
\begin{center}
\begin{tabular}{lccc}
\toprule
State & Target MM & V4 MM Achieved & MM District Minority \% \\
\midrule
Alabama      & 2 & \textbf{2} & 51.1\%, 50.5\% \\
Georgia      & 5 & \textbf{7} & 52.3\%--71.4\% \\
Louisiana    & 2 & 1           & 60.2\% \\
Mississippi  & 2 & 1           & 55.8\% (D1 at 49.99\%\dag) \\
South Carolina & 3 & 1         & 52.6\% \\
\bottomrule
\end{tabular}
\end{center}
{\small \dag Mississippi District 1 reaches 49.99\% minority---a stochastic boundary condition. Across repeated METIS runs with different random seeds, this district fluctuates across the 50\% threshold. The adaptive formula's effective $\alpha \approx 6.7\times$ is between the ablation's 5$\times$ and 10$\times$ configs; achieving both MM districts reproducibly would require the 10$\times$--50$\times$ range or a lower threshold.}

\textbf{Discrepancy from ablation headline}: The ablation study reports 80\% success (4/5 states) using per-state-optimal fixed configurations. The V4 production adaptive formula achieves 40\% of full legal targets (Alabama and Georgia only), with Louisiana and Mississippi each at 1/2 districts. Section~\ref{sec:discussion_adaptive} discusses this gap.

\subsection{Implementation}

\textbf{Partitioner}: METIS 5.1.0 (\texttt{gpmetis}) with flags \texttt{-contig} (enforce contiguity), \texttt{-minconn} (minimize subdomain connectivity), \texttt{-niter=100} (refinement iterations).

\textbf{Platform}: \texttt{redist} Rust binary with \texttt{--partition-mode metis-vra} (vra-aligned edge weights).

\textbf{Analysis}: Count MM districts with $\geq 50\%$ minority population. Report maximum minority concentration achieved and number of districts meeting threshold. For edge-weighted approach, additionally measure edge cut (unweighted and weighted) to assess compactness tradeoff.

All code and data available at [repository URL upon acceptance].
